\documentclass{article}\usepackage[margin=0.8in]{geometry}\begin{document}\section*{Signed graph forest closure}
Ten main declarations and nine new controls passed fresh replay, with nineteen exact standard-axiom readbacks. Eight modules returned native/gate exit zero. A separate false equality returned native/gate exit one.
\paragraph{Constructed forest witness.} Actual parent reachability identifies the child component. Add rank(parent)+1 to precisely that component. Every unchanged edge retains its decrease; the new edge strictly decreases to the surviving root. The proof derives forest preservation, child attachment and transfer of all old child-component paths to the surviving root.
\paragraph{Native edge composition.} The first returned root survives the second find. The composed theorem derives native reduced coefficients from the original edge and preserves Forest and all vertex/anchor equations for either union orientation. A later find has a finite successful execution. Assumptions include a valid initial forest/state, a true edge, distinct returned roots and unit-square decoded signs.
\paragraph{Controls.} Explicit signed-descendant finds feed both orientations. Root-to-descendant attachment forms a cycle under every rank. A nonroot child can safely be cut and attached to a distinct root, and this stronger accepted case is distinguished from unsafe nonroot-parent attachment.
\paragraph{Remaining interfaces.} Native size metadata, allocation, word/sign and array/counter bounds, compiled equivalence and the full same-root edge branch remain separate. The rank is a mathematical witness, not a native size update or complexity bound.
\paragraph{Custody.} Source-first plan, exact native binding, fresh invocation directories and retained failures are included. Independent review follows publication. Cached toolchain dependencies are trusted. No scientific execution or repository change occurred.
\end{document}